Sparse autoencoders (\sae{}s) trained on protein language model (\plm) activations decompose each layer into thousands of latent features, of which roughly 80\% do not align with any \swissprot concept. These \emph{dark features} have been hypothesized to encode biological mechanisms not yet catalogued by experimentalists, but no prior work tests this claim against causal, functional, or literature-grounded criteria at once. We audit the top dark features of the pretrained \interplm \sae for \esm (layer 4, 10{,}240 latents) on 1{,}500 human \swissprot proteins (422K residues) along three axes: (i) statistical structure against a random-residue null, (ii) causal influence on \esm's masked-LM logits via forward-hook ablation, and (iii) post-hoc annotation by \gptfive prompted on 15 maximally-activating sequence windows per feature. Eight of the top ten dark features detect the \tgekp inter-finger linker of C2H2 zinc fingers, and the remaining two detect the \dry and \npxxy motifs of class-A GPCRs; adding ``Zinc finger'' to the concept list lifts seven of those eight to F1 $\geq 0.6$. Single-feature ablation, however, fails for \emph{all} 30 top dark features under Benjamini--Hochberg FDR control: KL divergence at activating residues is on average $2^{-3.86}\times$ \emph{smaller} than at matched control residues. Our results refine the prevailing ``dark feature $=$ unknown biology'' framing: in \esm, top dark features encode well-known sub-domain motifs that \swissprot describes at coarser granularity, and high feature--concept correlation does not imply causal mechanistic role under single-feature ablation. We release the pipeline (\sae extraction $\rightarrow$ dark filter $\rightarrow$ \plm-grounded annotation) as a reusable audit tool for future \plm interpretability work.
